% Generated draft: requires mathematical and source review.
\documentclass[11pt]{article}
\usepackage[margin=25mm]{geometry}
\usepackage{fontspec}
\setmainfont{Latin Modern Roman}
\usepackage{amsmath,amssymb,mathtools,unicode-math}
\setmathfont{Latin Modern Math}
\usepackage{enumitem,needspace,xurl,hyperref}
\setlength{\parindent}{0pt}
\setlength{\parskip}{5pt}
\setlength{\emergencystretch}{3em}
\newcommand{\N}{\mathbb{N}}
\newcommand{\Z}{\mathbb{Z}}
\newcommand{\Q}{\mathbb{Q}}
\newcommand{\R}{\mathbb{R}}
\newcommand{\C}{\mathbb{C}}
\newcommand{\F}{\mathbb{F}}
\newcommand{\A}{\mathbb{A}}
\newcommand{\PP}{\mathbb{P}}
\newcommand{\E}{\mathbb{E}}
\newcommand{\eps}{\varepsilon}
\newcommand{\dd}{\,\mathrm{d}}
\begin{document}
\section*{UnsolvedMath ID: 700971}
\textbf{Source problem code:} catalogue entry
\subsection{Definitions and mathematical statement}
For a group \(G\), write \(\langle S\rangle\) for the subgroup generated by a subset \(S\subseteq G\), and use the convention \(b^g=g^{-1}bg\). An element \(x\in G\) has finite order if \(x^n=1_G\) for some integer \(n\geq 1\).

Is the following assertion true? For every infinite group \(G\) and every pair \(a,b\in G\) of finite-order elements satisfying \(G=\langle a,b\rangle\), the set
\[
\bigl\{g\in G : \langle a,g^{-1}bg\rangle\text{ is infinite}\bigr\}
\]
is infinite.

The assertion concerns infinitely many distinct elements \(g\), without requiring the corresponding conjugates or subgroups to be distinct.
\subsection{Short English statement}
If an infinite group is generated by two elements of finite order, must infinitely many conjugating elements produce an infinite subgroup generated by the first generator and the conjugate of the second?
\subsection{Sources}
\begin{itemize}
\item Local UnsolvedMath registry, permanent ID \(700971\), exact source code \texttt{catalogue entry}, titled \emph{Kourovka Notebook Problem 13.53}. The supplied record labels the problem open; this status is not independently verified here.
\item \emph{The Kourovka Notebook (Unsolved Problems in Group Theory)}, 13th issue (1995), Problem 13.53; proposed by A. I. Sozutov. The supplied background also identifies the 21st issue edition (2026), \texttt{arXiv:1401.0300v48}. The supplied collection description names E. I. Khukhro and V. D. Mazurov as editors.
\end{itemize}
\end{document}
